\section{Problem Formulation}
\label{sec:problem-formulation}

\subsection{Motivation}

Recognizing the category of a fashion item from an image is relevant to organizing product catalogs and supporting product search. This report studies a simplified version of this task using Fashion-MNIST. The objective is to compare how different neural network architectures learn from the same images within a reproducible experimental pipeline.

The comparison covers five model families: a linear classifier, a multilayer perceptron (MLP), a convolutional neural network (CNN), a long short-term memory network (LSTM), and a Transformer encoder. These models process the image in different ways: flattened pixels, a two-dimensional spatial grid, or a sequence of rows, columns, or patches. This makes the task suitable for examining how input representation and architectural assumptions relate to classification performance and computational cost.

The study is limited to the Fashion-MNIST benchmark. Its results should not be interpreted as evidence of performance on full-resolution product photographs, cluttered scenes, or production catalog systems.

\subsection{Input and Output}

Each sample is a grayscale image of size $28\times28$ pixels paired with one of ten class labels. The prediction unit is one image, and the task is single-label multiclass classification. After preprocessing, a batch of $B$ images is represented by a tensor $X\in\mathbb{R}^{B\times1\times28\times28}$.

For an individual image $x$, a model with parameters $\theta$ produces ten logits:
\begin{equation}
    f_{\theta}:\mathbb{R}^{1\times28\times28}\rightarrow\mathbb{R}^{10},
    \qquad z=f_{\theta}(x).
\end{equation}
The predicted class is
\begin{equation}
    \hat{y}=\operatorname*{arg\,max}_{k\in\{0,\ldots,9\}} z_k.
\end{equation}
All models return logits and are trained using cross-entropy loss. Softmax is not applied before the loss. Although their internal input representations differ, every model follows the same external input and output interface.

The main difficulties are distinguishing visually similar categories from low-resolution images and learning representations that generalize beyond the training set. The evaluation therefore considers both aggregate scores and the classes and examples on which each model makes errors.

\subsection{Research Questions}

The main research question is: \emph{How do the five architecture families compare on Fashion-MNIST in terms of predictive quality and computational cost under a shared data split and evaluation protocol?}

We examine the following supporting questions:
\begin{enumerate}
    \item How much do nonlinear models improve over the linear baseline in accuracy and macro-F1?
    \item How do spatial processing in the CNN and sequence processing in the LSTM and Transformer relate to performance and class confusions?
    \item How does representing an image as rows, columns, or patches affect the LSTM in the available experiments?
    \item What trade-offs arise between accuracy, macro-F1, parameter count, training time, and inference time?
\end{enumerate}

All main models use the same train, validation, and test partitions and evaluation metrics. Differences in model capacity, regularization, and training behavior are documented when interpreting the results; a performance difference alone does not isolate the causal effect of architecture. Numerical answers and supporting examples are presented in the results and discussion sections.
